\section{20261012}

If we have a line bundle with a section, we can describe its sections as rational functions. They correspond with regular sections and add the divisor of the distinguished section, we get an effective divisor.\par
This haunts algebraic geometers for half of their career. The choice of \(s\) induces a map:

\[
\Set{\text{sections of }L}\Leftrightarrow\Set{\text{rational functions on }X}
\]

In particular


\[
\Set{\text{regular sections}}\xrightarrow{\isom}\Set{\vf\text{ s.t. }\div\vf+\div(s)\geq 0}
\]

\red{Can we make this a diagram such that the bottom row }
\subsection{In toric land}

Suppose we have a torus invariant divisor \(D\). ¿What are the monomial rational functions describing the regular sections of the corresponding line bundle? Via this isomorphism? We want to impose the condition \(\div\vf+\div(s)\geq 0\) and recall our monomial rational functions are \(m\in M\). \par
The divisor of such an \(m\) is a sum over rays. We want: 
\[
\sum_{\rho\in\Sigma(1)}\braket{m}{u_\rho}D_\rho+a_\rho D_\rho\geq 0.
\]
So in particular for every ray: \(\braket{m}{u_\rho}\geq -a_\rho\). A combinatorialist looks at this and looks at it as a half space. Bunch of \(\rho\)'s gives intersections of half spaces and thus we get a polyhedron.\par

In conclusion, if \(D=\sum a_\rho D_\rho\) is a Cartier divisor on \(X_\Sg\), there is a canonical bijection between \(T\)-equivariant sections of \(L\) and the characters corresponding to lattice points in the polyhedron defined as 

\[
   \bigcap_{\rho\in\Sg(1)}\set{\braket{m}{u_\rho}\geq-a_\rho}. 
\]

\begin{Ex}
In \(\bP^2\), the fan is \(N_\bR\) and \(D=\set{Z=0}\). In \(M_\bR\) we get
\[
\begin{cases}
    \braket{(x,y)}{(1,0)}\geq 0\iff x\geq 0\\
    \braket{(x,y)}{(0,1)}\geq 0\iff y\geq 0\\
    \braket{(x,y)}{(-1,-1)}\geq 0\iff x+y\leq 1
\end{cases}
\]
We want the oreders of vanishing to be the right ones, so if the poly 
\end{Ex}